\section{Phonon and charge transport Monte Carlo}\label{app:E.12}

A \SI{100}{eV} deposit in a silicon crystal ends, after the cascades of \cref{chap:18,chap:19}, as some $10^4$ acoustic phonons of a few meV and, in a biased detector, a few tens of electrons and holes drifting toward the electrodes. The sensors on the surface see neither the deposit nor the cascade. They see phonons that arrive over microseconds after scattering off isotopes, decaying, bouncing off surfaces and crossing into metal films, and charges that arrive after drifting, emitting phonons of their own, and sometimes getting trapped. \Cref{chap:20} derives the physics of each step; combined in one simulation of a whole detector, they turn an event into the arrival distributions called collection kernels. The method is Monte Carlo solution of the Boltzmann equation for phonons and carriers, and the silicon case is worked with the rate constants and code conventions of the public detector-physics toolkit G4CMP~\cite{kelsey2023g4cmp,leman2012invited}.

\subsection{Computed quantities}

The state of a phonon is its wavevector $\vect{Q}$, its branch $\nu$ and its position; the state of a carrier is its crystal momentum, its valley or band, and its position. Between events each particle moves on a straight line with its group velocity (a carrier also accelerates in the applied field). Events are drawn at random with the rates of the scattering processes, and each event replaces the particle by its products. With many particles, the procedure samples the solution of the linear Boltzmann equation
\begin{equation}
  \partial_t f_i + \vect{v}_i\cdot\nabla_{\vect{r}} f_i + \vect{F}_i\cdot\nabla_{\vect{k}} f_i = \Bigl(\frac{\partial f_i}{\partial t}\Bigr)_{\mathrm{coll}},
  \label{eq:E.12-boltzmann}
\end{equation}
for the occupation $f_i(\vect{r}, \vect{k}, t)$ of each species $i$, with $\vect{v}_i$ its group velocity, $\vect{F}_i$ the force on it ($\pm e\vect{\mathcal{E}}$ for carriers in a field $\vect{\mathcal{E}}$, zero for phonons) and a collision term built from golden-rule rates~\cite{jacoboni1983monte}. The equation is linear because the athermal populations are dilute: phonons scatter off the static disorder and decay spontaneously, carriers emit into an empty phonon bath, and no process depends on the occupation of the other particles. The outputs are tallies at the boundaries: for every sensor, the energy and time of each phonon it absorbs, and for every electrode, the charge induced on it.

\subsection{Place on the hierarchy}

Every rate in the collision term is a lowest-order golden-rule rate. Isotope scattering is the second-order Born rate of a phonon off the random mass defects of \cref{sec:11.6}. Anharmonic decay is the rate of the cubic term of the lattice potential (\cref{sec:5.7}), at the lowest order at which it does not vanish, the cubic vertex squared. Neganov--Trofimov--Luke emission and intervalley scattering are one-phonon emission by a carrier through a deformation potential, $n_g = 2$ in the notation of the truncation label $\mathfrak{T}$ of \cref{sec:8.5}, with $n_B = 0$. The transport itself is incoherent: each particle is a wavepacket with a well-defined momentum, and interference between successive events is dropped. This holds when the scattering rate is small compared with the particle's frequency. For a silicon phonon at \SI{1}{THz} the isotope rate is \SI{2.4e6}{s^{-1}} (\cref{eq:E.12-iso} below) against $2\pi\times\SI{e12}{s^{-1}}$; even at the top of the acoustic band, \SI{15}{THz}, the ratio is $\sim10^{-3}$. It also needs mean free paths long compared with the wavelength, which they are by a factor $4\times10^5$ at \SI{1}{THz} (a wavelength of \SI{5.4}{nm} against a mean free path of \SI{2.2}{mm}). The Boltzmann description is therefore safe for the bulk; its approximations are concentrated in the rates and in the surface models, which is where the pitfalls below lie.

\subsection{Phonon rates}

Isotope scattering and anharmonic decay remove a phonon from its ballistic path. Isotope scattering redirects it, and may change its branch, at the rate~\cite{tamura1983isotope}, \cite[Eq.~(6)]{leman2012invited}
\begin{equation}
  \varGamma_{\mathrm{iso}} = B_{\mathrm{iso}}\left(\frac{E_{\mathrm{ath}}}{h}\right)^4 , \qquad
  B_{\mathrm{iso}} = 4\pi^3\,g_{\mathrm{md}}\,\Omega_{\mathrm{at}}\,\langle c^{-3}\rangle,
  \label{eq:E.12-iso}
\end{equation}
where $E_{\mathrm{ath}}/h$ is the phonon frequency, $\Omega_{\mathrm{at}}$ the volume per atom, $\langle c^{-3}\rangle = \tfrac13\sum_\nu c_\nu^{-3}$ the average of the inverse cubed sound speed over the three acoustic branches $\nu$ (the source writes $1/\langle c^3\rangle$ with the polarization-averaged cubed speed; the two agree when the speeds are equal, and $\langle c^{-3}\rangle$ is the average the Debye density of states contains, \cref{sec:20.2}), and $g_{\mathrm{md}} = \langle (M - \bar M)^2\rangle/\bar M^2$ the mass variance of the isotope mixture over the squared mean mass~\cite[Eqs.~(3)--(5)]{hernandez2025modeling}. It is the Rayleigh law of a point scatterer, valid as long as the wavelength is long compared with the interatomic spacing. For natural silicon, $g_{\mathrm{md}} = \num{2.02e-4}$, $\Omega_{\mathrm{at}} = \Omega/2 = \SI{20.0}{\angstrom^3}$ and, with $c_{\mathrm{LA}} = \SI{9.0}{km/s}$ and $c_{\mathrm{T}} = \SI{5.4}{km/s}$, $\langle c^{-3}\rangle = 1/(\SI{2.13e11}{m^3.s^{-3}}) = \SI{4.69e-12}{s^3.m^{-3}}$, and \cref{eq:E.12-iso} gives $B_{\mathrm{iso}} = 124.0\times\num{2.02e-4}\times\SI{2.00e-29}{m^3}\times\SI{4.69e-12}{s^3.m^{-3}} = \SI{2.35e-42}{s^3}$, in line with the tabulated $\SI{2.33e-42}{s^3}$ and with the value $\SI{2.43e-42}{s^3}$ that most transport codes use~\cite[Table~4]{hernandez2025modeling}, \cite[Table~I]{leman2012invited}. The direction and branch of the scattered phonon are then drawn with the weight $|\hat{\vect{e}}\cdot\hat{\vect{e}}'|^2$ of the polarization vectors of the initial and final modes, divided by the cube of the final frequency~\cite[Eq.~(7)]{leman2012invited}. Repeated isotope scattering drives the branch populations of silicon to slow transverse \SI{55}{\percent}, fast transverse \SI{35}{\percent} and longitudinal \SI{10}{\percent}~\cite[Sect.~III\,E]{leman2012invited}.

Anharmonic decay splits a longitudinal phonon into two of lower energy, at the rate~\cite{tamura1985spontaneous}, \cite[Eq.~(8)]{leman2012invited}
\begin{equation}
  \varGamma_{\mathrm{anh}} = A_{\mathrm{anh}}\left(\frac{E_{\mathrm{ath}}}{h}\right)^5 ,
  \label{eq:E.12-anh}
\end{equation}
with $A_{\mathrm{anh}}$ fixed by the second- and third-order elastic constants in the isotropic continuum~\cite[Eq.~(8)]{leman2012invited}. Transverse phonons hardly decay at low frequency. The two channels, $\mathrm{L}\to\mathrm{L}+\mathrm{T}$ and $\mathrm{L}\to\mathrm{T}+\mathrm{T}$, have energy distributions of the daughters that the isotropic approximation fixes in closed form. For $\mathrm{L}\to\mathrm{L}'+\mathrm{T}'$, with $x = E_{\mathrm{L}'}/E_{\mathrm{L}}$ and $\delta = c_{\mathrm{LA}}/c_{\mathrm{T}}$, the rate is distributed as $x^{-2}(1 - x^2)^2[(1 + x)^2 - \delta^2(1 - x)^2][1 + x^2 - \delta^2(1 - x)^2]^2$ over $(\delta - 1)/(\delta + 1) < x < 1$, and the $\mathrm{T}+\mathrm{T}$ channel has a polynomial distribution of the same kind with coefficients set by the elastic constants~\cite[Eqs.~(9) and~(12)]{leman2012invited}. The daughter directions then follow from energy and momentum conservation; for the longitudinal daughter, $\cos\theta_{\mathrm{L}'} = [1 + x^2 - \delta^2(1 - x)^2]/(2x)$~\cite[Eq.~(10)]{leman2012invited}. For silicon, Tamura's isotropic calculation gives $A_{\mathrm{anh}} = \SI{7.41e-56}{s^4}$, as tabulated by Leman~\cite[Table~I]{leman2012invited}; a recent recalculation from cryogenic elastic constants gives $\SI{1.15e-55}{s^4}$~\cite[Table~3]{hernandez2025modeling}. No direct measurement exists for silicon, so $A_{\mathrm{anh}}$ is uncertain by at least the factor 1.5 between the two.

\Cref{tab:E-phrates} evaluates \cref{eq:E.12-iso,eq:E.12-anh} for silicon. The ratio of the two rates, $\varGamma_{\mathrm{anh}}/\varGamma_{\mathrm{iso}} = (A_{\mathrm{anh}}/B_{\mathrm{iso}})(E_{\mathrm{ath}}/h)$, equals $(E_{\mathrm{ath}}/h)/\SI{33}{THz}$ and is below one everywhere in the acoustic band. Decay is nevertheless what lowers the frequency, since isotope scattering conserves energy. A \SI{10}{THz} longitudinal phonon decays within \SI{0.13}{ns}. Its daughters near \SI{3}{THz} live some \SI{55}{ns} before decaying again, scattering off isotopes about ten times in between, and below \SI{1}{THz} the decay time exceeds \SI{10}{\micro s}, longer than the time a phonon needs to reach a surface of a centimeter-sized crystal.

\begin{table}[t]
  \caption{Bulk phonon rates in natural silicon from \cref{eq:E.12-iso,eq:E.12-anh} with $B_{\mathrm{iso}} = \SI{2.43e-42}{s^3}$ and $A_{\mathrm{anh}} = \SI{7.41e-56}{s^4}$~\cite{leman2012invited}. The anharmonic rate applies to longitudinal phonons. The isotope mean free path uses the transverse speed \SI{5.4}{km/s}~\cite{leman2012invited}.}
  \label{tab:E-phrates}
  \centering
  \begin{tabular}{ccccc}
    \toprule
    $E_{\mathrm{ath}}/h$ (THz) & $E_{\mathrm{ath}}$ (meV) & $\varGamma_{\mathrm{iso}}$ (s$^{-1}$) & isotope mean free path & $\varGamma_{\mathrm{anh}}$ (s$^{-1}$) \\
    \midrule
    0.5 & 2.07 & \num{1.5e5} & \SI{3.6}{cm} & \num{2.3e3} \\
    1 & 4.14 & \num{2.4e6} & \SI{2.2}{mm} & \num{7.4e4} \\
    3 & 12.4 & \num{2.0e8} & \SI{27}{\micro m} & \num{1.8e7} \\
    10 & 41.4 & \num{2.4e10} & \SI{0.22}{\micro m} & \num{7.4e9} \\
    \bottomrule
  \end{tabular}
\end{table}

\subsection{Free flight, focusing and surfaces}

Between events a phonon moves with the group velocity of its branch, obtained from the Christoffel equation of continuum elasticity, $\rho_T\omega^2\hat e_i = C_{ijkl}Q_jQ_k\hat e_l$, for the unit polarization $\hat{\vect{e}}$ of the elastic wave and the elastic tensor $C_{ijkl}$ (\cref{app:C.18}). In a cubic crystal the group velocity is not parallel to $\vect{Q}$, and an isotropic distribution of wavevectors becomes a strongly peaked distribution of energy flux, the phonon focusing that concentrates ballistic phonons along certain crystal directions~\cite[Sect.~2.4]{hernandez2025modeling}. Using the continuum dispersion is an approximation for the phonons near the top of the acoustic band, whose dispersion bends over; it is good for the few-THz phonons that dominate after the first decays.

A phonon that reaches a surface is reflected specularly, reflected diffusely, down-converted into lower-energy phonons, or transmitted into whatever is attached. The bare crystal surface is described by the probabilities of these outcomes, which are inputs of the simulation, fitted to measured pulse shapes for a given surface preparation. An interface to a sensor film is described by a transmission probability, taken from the acoustic mismatch model or the diffuse mismatch model~\cite{swartz1989thermal}. \Cref{sec:20.3} derives both and the reciprocity they obey. Measured silicon-to-aluminum values, $\bar t_{12}P_{\mathrm{abs}} = 0.3$--0.55, are absorption probabilities and not transmissions~\cite{martinez2019measurements}, and a simulation that uses them must use them in that role.

\subsection{Films}

Inside a superconducting film the phonon breaks Cooper pairs, the resulting quasiparticles relax by emitting phonons, and phonons with energy above $2\Delta_{\mathrm{BCS}}$ break further pairs; this is the Kaplan cascade of \cref{sec:20.4}~\cite{kaplan1976quasiparticle}. A full simulation tracks quasiparticles and phonons in the film. The lumped model used in detector codes treats the film as a single object with a thickness $d_{\mathrm{f}}$, a gap $\Delta_{\mathrm{BCS}}$ and an energy-dependent phonon mean free path $\mathfrak{l}_{\mathrm{f}}(E_{\mathrm{ath}})$. An incident phonon escapes back into the substrate unabsorbed with probability $\exp[-4d_{\mathrm{f}}/\mathfrak{l}_{\mathrm{f}}(E_{\mathrm{ath}})]$, the factor 4 being twice the thickness, for the round trip, times two from the average over incidence angles~\cite[Eq.~(17)]{kelsey2023g4cmp}. Otherwise its energy goes to two quasiparticles, which share it with the BCS density of states weight; quasiparticles below $3\Delta_{\mathrm{BCS}}$ are counted as collected, and those above emit a phonon and continue~\cite[Sect.~2.4]{kelsey2023g4cmp}. Phonons below $2\Delta_{\mathrm{BCS}}$ cannot break pairs and are either absorbed directly with a set probability or returned to the substrate. The output is the quasiparticle yield $Y_{\mathrm{qp}}(E_{\mathrm{ath}})$, sampled event by event, and the phonons the film sends back into the crystal.

\subsection{Carrier transport}

In a biased detector the electrons and holes of the cascade drift in the field $\vect{\mathcal{E}}$, computed beforehand for the electrode geometry by a finite-element solution of the Laplace equation. A carrier accelerates until it moves faster than the longitudinal sound speed, and then emits acoustic phonons, the acoustic analogue of Cherenkov radiation (\cref{sec:20.1}). These are the Neganov--Trofimov--Luke phonons, and they carry the work $eV_{\mathrm{bias}}$ that the bias does on each pair, the Luke gain of \cref{sec:19.5}~\cite{luke1988voltage}. For a carrier with an isotropic mass $m$ and momentum $k$, and the acoustic deformation potential $D^n$ of its band (\cref{sec:6.2}) taken isotropic, the golden-rule emission rate is $\varGamma_{\mathrm{L}}(k)$ of \cref{eq:20.1-luke}, with $k_{\mathrm{s}} = m\,c_{\mathrm{LA}}$ the momentum of a carrier moving at the sound speed. The polar angle $\theta$ of the emitted phonon, measured from $\vect{k}$, is drawn from $P(\theta) \propto (\cos\theta - k_{\mathrm{s}}/k)^2\sin\theta$ for $0 \le \theta \le \arccos(k_{\mathrm{s}}/k)$~\cite[Eq.~(9)]{kelsey2023g4cmp}. The rate vanishes below $k_{\mathrm{s}}$ and grows as $k^2$ at large $k$. It counts longitudinal phonons only; a carrier faster than the transverse sound speed, which is lower, can also emit transverse ones, and a model that omits them shifts the Luke phonons toward the longitudinal branch.

Holes in silicon and germanium are close to isotropic, and \cref{eq:20.1-luke} applies directly. Electrons sit in anisotropic valleys, six along $\langle100\rangle$ in silicon and four along $\langle111\rangle$ in germanium, each with the mass tensor of components $m^*_\parallel$ along the valley axis and $m^*_\perp$ across it (\cref{sec:3.3}). A Herring--Vogt transformation rescales the momentum along each principal axis by $(m_{\mathrm{c}}/m^*_i)^{1/2}$, with the conductivity mass $3/m_{\mathrm{c}} = 1/m^*_\parallel + 2/m^*_\perp$, and turns the valley into an isotropic one. The field and the emission are handled there as for holes, and the results are transformed back~\cite{herring1956transport}, \cite[Eqs.~(11)--(12)]{kelsey2023g4cmp}. The transformation is not unitary, and the emitted phonon's momentum is transformed back by keeping its magnitude, which conserves energy, and transforming only its direction~\cite[Sect.~V\,C]{leman2012invited}. Electrons also scatter between valleys by emitting or absorbing intervalley phonons, which turns their trajectories toward the field direction. The intervalley rates are parametrized, from measurements in silicon and from a model in germanium, rather than computed from first principles~\cite[Sect.~2.3]{kelsey2023g4cmp}; the drift speed against field is then a check of the whole carrier model, which reproduces measurements in germanium~\cite[Fig.~7]{kelsey2023g4cmp}.

Impurities enter as two processes, each with a mean free path that is fitted to the device. One is trapping, in which a carrier is captured by a neutral impurity and stops, and the other is impact ionization of an impurity, in which a carrier frees a second charge from a charged one~\cite[Eqs.~(13)--(16)]{kelsey2023g4cmp}. Trapping reduces the collected charge and the Luke phonons, impact ionization increases them, and both fill the regions between the single-pair peaks of a charge-resolving silicon detector, from whose shape the four or six mean free paths are fitted~\cite[Fig.~6]{kelsey2023g4cmp}. In a uniform field with trapping alone, the simulation must reproduce the Hecht relation of \cref{sec:20.1}; that is the first check of a trapping model.

The step size for carriers is limited by the requirement that the rate stays nearly constant during a step. Taking a step of half the current mean free time conserves energy, in the sense that the energy of the Luke phonons emitted equals the work done by the field, and longer steps produce too few Luke phonons~\cite[Sect.~V\,D]{leman2012invited}. The number of Luke phonons is large: each pair crossing a bias $V_{\mathrm{bias}}$ releases $eV_{\mathrm{bias}}$, and a code that assumes \SI{2}{meV} per Luke phonon in estimating its event size~\cite[Sect.~3.4]{kelsey2023g4cmp} expects $\SI{100}{eV}/\SI{2}{meV} = 5\times10^4$ phonons per pair at \SI{100}{V}.

\subsection{Collection kernels}

The collection kernel of \cref{sec:20.5} is the probability that energy deposited at $\vect{r}_{\mathrm{ev}}$, in phonons of energy $E_{\mathrm{ath}}$, is absorbed in sensor $s$ within a time window. It is tallied by launching phonons of each energy from a grid of event positions and counting where they end. For charge, the analogous kernel is the induced charge on each electrode, computed from the trajectories with Ramo's theorem and the weighting potentials $\varphi_s$ of \cref{sec:20.1}. Kernels separate the transport from the event: once tabulated, they are convolved with the deposit and the phonon spectrum of any event, from the created-phonon spectra of \cref{chap:19} to the down-converted spectra of \cref{sec:20.2}, without rerunning the transport.

The statistical error of a kernel bin with $n$ absorbed out of $N_{\mathrm{l}}$ launched phonons is binomial, $[n(1 - n/N_{\mathrm{l}})]^{1/2}/N_{\mathrm{l}}$. For the collection contrasts of \cref{sec:20.5}, which are differences between kernels of similar size, the useful technique is to run the two cases with common random numbers, so that their statistical fluctuations are correlated and cancel in the difference.

\subsection{Validation and the closure checks}

Four checks cost nothing: energy, charge, rate and step closure. For energy closure, the deposit, plus the work of the bias on the carriers, must equal the energy absorbed in the sensors, plus the energy lost at bare surfaces, plus the energy in trapped carriers and in particles still in flight at the end, to the round-off of the tallies. A missing term in this balance signals a process that loses energy without booking it, which is the most common silent bug in such codes. For charge closure, the charge induced on all electrodes, summed with the trapped charges weighted by their weighting potentials, must equal the charge created. For rate closure, in a large block with surfaces removed and one process switched on at a time, the measured mean free paths must follow \cref{eq:E.12-iso,eq:E.12-anh} with the input constants, with the slopes $E_{\mathrm{ath}}^{-4}$ and $E_{\mathrm{ath}}^{-5}$; this has been done for the materials of the recent rate tabulation~\cite[Sect.~3]{hernandez2025modeling}. For step closure, halving every time step must leave the kernels unchanged within their statistical errors.

Beyond these internal checks, a run must report its agreement with external benchmarks that test the physics rather than the bookkeeping. They are carrier drift velocities against field, the Hecht curve for a known trap density, the ballistic focusing pattern of a thin crystal, the quasiparticle yield $Y_{\mathrm{qp}}(E_{\mathrm{ath}})$ against the Kaplan cascade of \cref{sec:20.4}, and, where it exists, a measured pulse shape.

\subsection{Cost}

The cost is set by the number of particle steps. A \SI{10}{keV} deposit in silicon that has down-converted to \SI{1}{THz} phonons ($\SI{4.1}{meV}$) holds $2.4\times10^6$ phonons. In a crystal a centimeter across, each phonon scatters off isotopes about every \SI{2}{mm} (\cref{tab:E-phrates}) and bounces off the surfaces many times before a sensor covering a few percent of the surface absorbs it. At a hundred steps per phonon the event needs $\sim10^8$--$10^9$ steps, and a biased detector adds the $5\times10^4$ Luke phonons per pair per \SI{100}{V}. With full tracking, a few-keV event takes a day or more of central processing unit (CPU) time, and with the down-sampling described next about an hour per event~\cite[Sect.~3.4 and Fig.~8]{kelsey2023g4cmp}. Collection kernels are much cheaper than events, since each launched phonon is an independent sample and no event needs to be complete.

Down-sampling generates only a fraction $f$ of the phonons and carriers and gives each the weight $1/f$~\cite[Sect.~3.4]{kelsey2023g4cmp}. It conserves the mean of every tally. It does not conserve the fluctuations: a weighted event has a variance $1/f$ times too large for the collected energy, and a variance computed from weighted samples is not the physical variance. Kernels and means may be computed with down-sampling; the measured-charge Fano factor $F_{\mathrm{meas}}$ of \cref{sec:20.1}, and any statement about event-by-event resolution, may not.

\subsection{Pitfalls}

No closure check sees the errors in three inputs: the sound speed used to convert a decay calculation to the fifth-power law, the starting phonon spectrum, and the surface, intervalley and trapping parameters. The fifth-power law of \cref{eq:E.12-anh} amplifies any inconsistency in that sound speed, as the recalculation found for sapphire~\cite[Sect.~3.1]{hernandez2025modeling}. The transport is usually started from phonons at the Debye frequency, on the grounds that the first decays erase the details~\cite[Sect.~III\,B]{leman2012invited}, which is safe for kernels at long times but not for the prompt signal near the event, the subject of the memory-loss arguments of \cref{sec:20.2}. The surface parameters are fits to the device in which they were measured, so that transferring them to a crystal from a different boule is an assumption to be stated.

\subsection{Silicon}

\Cref{tab:E-transport-si} lists the silicon inputs of a phonon and charge transport run, with the checks that a run must report alongside its kernels.

\begin{table}[t]
  \caption{Phonon and charge transport in silicon: inputs, the quantities they control, and the checks a run must report. Rates are from \cref{eq:E.12-iso,eq:E.12-anh} with the cited constants.}
  \label{tab:E-transport-si}
  \centering
  \begin{tabular}{p{0.36\linewidth}p{0.52\linewidth}}
    \toprule
    Input or check & Silicon \\
    \midrule
    $B_{\mathrm{iso}}$; $A_{\mathrm{anh}}$ & \SI{2.43e-42}{s^3}; \SIrange{7.41e-56}{1.15e-55}{s^4}~\cite{leman2012invited,hernandez2025modeling} \\
    Sound speeds L, T & 9.0, \SI{5.4}{km/s}~\cite{leman2012invited} \\
    Isotope mean free path at \SI{1}{THz} & \SI{2.2}{mm} \\
    Equilibrium branch populations slow transverse (ST) : fast transverse (FT) : longitudinal & 55 : 35 : 10~\cite{leman2012invited} \\
    Si $\to$ Al absorption probability $\bar t_{12}P_{\mathrm{abs}}$ & 0.3--0.55~\cite{martinez2019measurements} \\
    Luke phonons per pair at \SI{100}{V} & $5\times10^4$ (at \SI{2}{meV} each) \\
    Carrier model & six $\langle100\rangle$ electron valleys, Herring--Vogt; isotropic holes; fitted intervalley, trapping and impact-ionization mean free paths \\
    Closure checks & energy, induced charge, rate slopes $E_{\mathrm{ath}}^{-4}$ and $E_{\mathrm{ath}}^{-5}$, time step \\
    Benchmarks & drift velocity, Hecht curve, focusing pattern, $Y_{\mathrm{qp}}(E_{\mathrm{ath}})$, pulse shape \\
    \bottomrule
  \end{tabular}
\end{table}
